\documentclass[11pt,a4paper]{article}
\usepackage[utf8]{inputenc}
\usepackage[T1]{fontenc}
\usepackage[portuguese]{babel}
\usepackage{geometry}
\geometry{margin=2.5cm}
\usepackage{listings}
\usepackage{xcolor}
\usepackage{amsmath}
\usepackage{amssymb}
\usepackage{hyperref}
\usepackage{enumitem}
\usepackage{tcolorbox}
\usepackage{tabularx}
\usepackage{booktabs}
\usepackage{graphicx}

\definecolor{codebg}{rgb}{0.95,0.95,0.95}
\definecolor{codegreen}{rgb}{0.1,0.5,0.1}
\definecolor{codepurple}{rgb}{0.5,0.1,0.5}
\definecolor{codegray}{rgb}{0.5,0.5,0.5}

\lstdefinestyle{python}{
    backgroundcolor=\color{codebg},
    commentstyle=\color{codegray}\itshape,
    keywordstyle=\color{codepurple}\bfseries,
    stringstyle=\color{codegreen},
    basicstyle=\ttfamily\small,
    breaklines=true,
    frame=single,
    language=Python,
    showstringspaces=false,
    tabsize=4,
    literate={á}{{\'a}}1 {ã}{{\~a}}1 {é}{{\'e}}1 {í}{{\'i}}1 {ó}{{\'o}}1 {õ}{{\~o}}1 {ú}{{\'u}}1 {ç}{{\c{c}}}1 {Á}{{\'A}}1 {Ã}{{\~A}}1 {É}{{\'E}}1 {Í}{{\'I}}1 {Ó}{{\'O}}1 {Õ}{{\~O}}1 {Ú}{{\'U}}1 {Ç}{{\c{C}}}1 {â}{{\^a}}1 {ê}{{\^e}}1 {ô}{{\^o}}1 {û}{{\^u}}1 {Â}{{\^A}}1 {Ê}{{\^E}}1 {Ô}{{\^O}}1 {Û}{{\^U}}1 {à}{{\`a}}1 {è}{{\`e}}1 {ì}{{\`i}}1 {ò}{{\`o}}1 {ù}{{\`u}}1 {ä}{{\"a}}1 {ö}{{\"o}}1 {Ä}{{\"A}}1 {Ö}{{\"O}}1
}
\lstset{style=python}

\newtcolorbox{objetivos}[1][]{
  colback=green!5!white,
  colframe=green!40!black,
  title=O que você vai aprender nesta página,
  fonttitle=\bfseries,
  #1
}

\newtcolorbox{exercicio}[1]{
  colback=blue!5!white,
  colframe=blue!40!black,
  title=#1,
  fonttitle=\bfseries
}

\newtcolorbox{solucao}{
  colback=yellow!5!white,
  colframe=orange!60!black,
  title=Solução proposta,
  fonttitle=\bfseries
}

\title{\textbf{Data Analysis with Python 2024--2025}\\Parte 7: Project Work\\\large Caderno de Orientações e Exercícios}
\author{Tradução e Adaptação: University of Helsinki --- MOOC.fi}
\date{}

\begin{document}

\maketitle

\section*{Nota Importante}
Este material é uma tradução e adaptação baseada no curso \textit{Data Analysis with Python 2024-2025}, oferecido pela \textbf{Universidade de Helsinque (MOOC.fi)}.

\tableofcontents
\newpage

\section{Caderno de Tarefas - Projeto Final (Chapter 7)}

\subsection*{Diretrizes Gerais para Submissão}

\begin{exercicio}{Orientações Práticas do Projeto}
\begin{itemize}
    \item \textbf{Escolha do Tema:} O aluno deve selecionar obrigatoriamente um dentre os três projetos disponíveis: \emph{Sequence Analysis}, \emph{Regression Analysis} ou \emph{Fossil Data Analysis}.
    \item \textbf{Estrutura do Arquivo:} O trabalho final deve ser obrigatoriamente salvo e entregue em formato de arquivo Jupyter Notebook com a extensão \texttt{.ipynb}.
    \item \textbf{Documentação Textual:} Nos projetos que exigem justificativas e explicações teóricas (como a análise de sequências), preencha de forma clara as caixas de texto contendo a descrição da lógica e a análise crítica dos resultados obtidos.
    \item \textbf{Integridade do Código:} Nos notebooks estruturados com marcações automáticas (ex: linhas \texttt{\# exercise x}), abstenha-se de apagar ou modificar estas marcações para evitar falhas nos testes automatizados do TMC.
\end{itemize}
\end{exercicio}

\subsection*{Atividades de Revisão por Pares (Peer Review)}

\begin{exercicio}{Validação de Trabalhos de Terceiros}
\begin{itemize}
    \item \textbf{Ambiente Isolado:} Sempre execute os testes dos colegas em diretórios temporários do sistema operacional para prevenir a perda ou sobrescrita acidental das suas próprias respostas já submetidas.
    \item \textbf{Geração do Link de Envio:} Utilize o comando \textit{Send Exercise to TMC Paste} no editor ou a ferramenta de linha de comando para gerar a URL de compartilhamento de código no formato \url{https://tmc.mooc.fi/paste/}.
    \item \textbf{Obrigação de Feedback:} Após submeter o seu próprio arquivo de projeto na plataforma oficial do MOOC, o sistema exigirá a emissão de avaliações para três outros relatórios de discentes da turma.
\end{itemize}
\end{exercicio}

\end{document}
